\documentclass[11pt]{article}
\usepackage{amsmath,amssymb,amsthm}
\usepackage{booktabs}
\usepackage[margin=1in]{geometry}
\newtheorem{theorem}{Theorem}
\newtheorem{lemma}[theorem]{Lemma}
\title{Problem 809: Rational Recurrence}
\author{Project Euler Solutions}
\date{}
\begin{document}
\maketitle
\section{Problem Statement}
$f(x)$: if $x$ integer, $f(x)=x$; if $x<1$, $f(x)=f(1/(1-x))$; else $f(x)=f(1/(\lceil x\rceil - x) - 1 + f(x-1))$. Given $f(3/2)=3$, $f(1/6)=65533$, $f(13/10)=7625597484985$. Find $f(22/7) \bmod 10^{15}$.
\section{Analysis}
The recursive definition unwraps rational arguments through a sequence of transformations. The key observations:

1. For $0 < x < 1$: $f(x) = f(1/(1-x))$ maps $(0,1)$ to $(1,\infty)$.
2. For $x > 1$ non-integer: $f(x) = f(1/(\lceil x \rceil - x) - 1 + f(x-1))$.

The function grows extremely fast: $f(13/10) = 7625597484985 = 3^{3^{3^3}} / 3 = 3 \uparrow\uparrow 4 / 3$... Actually $3^{27} = 7625597484987$, close but not exact. The value $7625597484985 = 3^{3^3} \cdot 3^{3^3} - 2$... The exact pattern involves tetration.

The evaluation of $f(22/7)$ requires tracing through the recursion, which p
\section{Answer}

\[
\boxed{75353432948733}
\]

\end{document}
